% One slide: the given feedback controller, and the five steps that put it in the training loop.
% Compile: pdflatex controller-slide.tex   (twice)
\documentclass[aspectratio=169,11pt]{beamer}

\usetheme{default}
\usecolortheme{seahorse}
\setbeamertemplate{navigation symbols}{}
\setbeamertemplate{itemize items}[circle]
\setbeamertemplate{enumerate items}[default]
\setbeamerfont{frametitle}{size=\large}
\setbeamercolor{block title}{bg=structure!20,fg=black}

\usepackage[T1]{fontenc}
\usepackage[utf8]{inputenc}
\usepackage{lmodern}
\usepackage{amsmath,amssymb,bm}
\usepackage{microtype}

\newcommand{\R}{\mathbb{R}}
\newcommand{\diag}{\operatorname{diag}}
\newcommand{\Cfb}{C_{\mathrm{fb}}}
\newcommand{\Cnorm}{C_{\mathrm{n}}}
\newcommand{\wb}{\omega_{\mathrm{b}}}
\newcommand{\udata}{u_{\mathrm{d}}}
\newcommand{\umodel}{u_{\mathrm{m}}}
\newcommand{\ydata}{y_{\mathrm{d}}}
\newcommand{\ymodel}{y_{\mathrm{m}}}
\newcommand{\Yop}{Y_{\mathrm{op}}}

\begin{document}

\begin{frame}{Closing the loop in training}

\begin{columns}[T,onlytextwidth]

%---------------------------------------------------------------- left
\begin{column}{0.43\textwidth}
\begin{block}{The controller}
\scriptsize
\begin{itemize}\setlength\itemsep{0.2em}
  \item Three independent SISO loops, one per stage axis $X_1, X_2, Y$, diagonal in the stage
        frame
  \item Loop shaping: integrator, lead-lag, first-order roll-off, $100$~Hz bandwidth
  \item Rule of thumb from the FP model: the shape fixes the phase margin, one gain per channel
        puts crossover at $100$~Hz
  \vspace{-0.2em}
  \item Tustin discretised, order 3 per channel, so 9 states
\end{itemize}
\vspace{0.2em}
{\scriptsize
\begin{align*}
  \Cnorm(s) &= \frac{s+\wb/6}{s}\cdot\frac{s+\wb/3}{s+3\wb}\cdot\frac{10\wb}{s+10\wb} \\[-0.2em]
  \kappa_j  &= \frac{1}{\bigl|\mathrm{sys}_{jj}(i\wb)\,\Cnorm(i\wb)\bigr|},
               \quad \wb = 2\pi f_{\mathrm{bw}}
\end{align*}}
\vspace{-1.2em}
\end{block}
\end{column}

%---------------------------------------------------------------- right
\begin{column}{0.55\textwidth}
\begin{block}{How it works in training}
\scriptsize
\begin{enumerate}\setlength\itemsep{0.25em}
  \item \textbf{Take it as a filter.} $(A_c,B_c,C_c,D_c)$, 9 states, fixed coefficients.
  \item \textbf{Wrap it around the model.} Error is reference minus model output, controller force
        injected at the model input.
  \item \textbf{Subtract the machine loop.} Reference and multisine are identical in both and
        cancel:
        \vspace{-0.3em}
        \[
          \umodel = \udata + \Cfb\,\bigl(\ydata - \ymodel\bigr)
        \]
        \vspace{-1.1em}
        \\ so training needs only the recorded input and output.
  \item \textbf{Fix the order.} Model output first.
  \item \textbf{Keep it separate.} Not part of the model state, so $n_x$ stays $8$, and $\Cfb$ is
        rebuilt per record at that record's $\Yop$.
\end{enumerate}
\end{block}
\end{column}

\end{columns}

\vspace{0.4em}
{\footnotesize
Only approximation: controller state reset to zero at each window start.}

\end{frame}

\end{document}
